\chapter{Large Deviations}
\begin{center}\itshape
The central limit theorem describes the usual fluctuations of an average. Insurers, engineers and physicists care about the unusual ones:\\
the year when claims are twice the average, the code word that the noise turns into another, the magnet that chooses a direction.\\
Their probabilities are exponentially small, and the exponent is an entropy.
\end{center}

An insurance company that collects premiums slightly above the expected claims will be ruined only if the claims of the year are far above their
mean: an event far in the tail of the sum. The Scandinavian actuaries of the early twentieth century needed its probability, and they found that the
Gaussian approximation is useless there. In 1932 Fredrik Esscher, a Swedish actuary, proposed a trick: change the probabilities, multiplying the law of each claim by
$e^{\theta x}$ and renormalizing, so that the rare event becomes typical, then pay for the change with an explicit factor. Harald Cramér, who worked on
insurance mathematics all his life, turned the trick into a theorem for a probability conference held in Geneva in 1937, published in 1938. It was the
first \emph{large deviation} theorem.

Twenty years later Ivan Sanov, in Moscow, asked the same question one level up: not ``how likely is a strange average?'' but ``how likely is a strange
histogram?'', and found that the answer is the relative entropy of Tome~3. In the 1960s S.~R.~Srinivasa Varadhan saw a single principle behind these
results and many others, from diffusions to statistical mechanics; the Abel Prize of 2007 was awarded to him ``for his fundamental contributions to
probability theory and in particular for creating a unified theory of large deviations''.

This chapter proves Cram\'er's theorem on the line and Sanov's theorem for finite alphabets, and draws four consequences: the shape of a rare event (the Gibbs
conditioning principle), the best possible statistical test (a lemma of Chernoff and Stein), the asymptotics of exponential integrals (Varadhan's lemma), and the
spontaneous magnetization of a toy magnet.

\begin{reminder}[from Tome 3]
For Bernoulli$(p)$ variables and $p<a<1$, $\frac1n\ln P(S_n\ge na)\to-I_p(a)$, where $I_p(a)=a\ln\frac ap+(1-a)\ln\frac{1-a}{1-p}$ (Tome~3, Theorem~4.7); the upper bound came
from Markov's inequality applied to $e^{\lambda S_n}$. The relative entropy $D(\nu\,\|\,\mu)=\sum_a\nu(a)\log\frac{\nu(a)}{\mu(a)}$ is non-negative and vanishes only for $\nu=\mu$
(Gibbs' inequality, Tome~3, Proposition~4.16). Tome~3 used $\log_2$; \emph{in this chapter all logarithms are natural}, and $D$ is measured in nats. The weak law of large
numbers holds for i.i.d.\ variables with finite variance (Tome~3, Theorem~4.2).
\end{reminder}

\section{Cram\'er's theorem}

Let $X$ be a real random variable with law $\mu$, not almost surely constant, and assume throughout this section that
\begin{equation}\label{eq:cramercond}
\Lambda(\theta)=\log\mathbb Ee^{\theta X}<\infty\qquad\text{for every }\theta\in\R.
\end{equation}
This holds for bounded, Gaussian and Poisson variables; Exercise~\ref{ex:heavyld} shows what goes wrong without it. Let $m=\mathbb EX$, and let $L<M$ be the essential
infimum and supremum of $X$ (possibly infinite). For each $\theta$, the \emph{tilted law} is $\mu_\theta(\dd x)=e^{\theta x-\Lambda(\theta)}\mu(\dd x)$: Esscher's trick.

\begin{lemma}[the cumulant generating function]\label{lem:Lambda}
$\Lambda$ is $C^\infty$, $\Lambda'(\theta)$ is the mean of $\mu_\theta$ and $\Lambda''(\theta)>0$ its variance. The map $\Lambda'$ is an increasing bijection from $\R$ onto $(L,M)$.
\end{lemma}
\begin{proof}
For $|\theta'-\theta|\le1$ and $k\ge0$, $|x|^ke^{\theta'x}\le C_k\bigl(e^{(\theta+2)x}+e^{(\theta-2)x}\bigr)$, an integrable bound, so we may differentiate $\mathbb Ee^{\theta X}$ under the
expectation as often as we like. Then $\Lambda'=\mathbb E[Xe^{\theta X}]/\mathbb Ee^{\theta X}$ is the mean of $\mu_\theta$, and $\Lambda''$ is its variance, positive since $\mu_\theta$ has the same null sets as
$\mu$ and is not a point mass. So $\Lambda'$ is continuous and strictly increasing, and its values lie in $(L,M)$, as means of laws carried by $[L,M]$ that are not point masses at
an endpoint.

It remains to see that $\Lambda'(\theta)\to M$ as $\theta\to\infty$ (the limit at $-\infty$ is symmetric). Let $b<b'<M$. Then
\[
\Lambda'(\theta)-b'=\frac{\mathbb E[(X-b')e^{\theta X}]}{\mathbb Ee^{\theta X}}\ge-\frac{\mathbb E[(b'-X)^+e^{\theta X}]}{\mathbb Ee^{\theta X}}.
\]
Split the numerator on the right. On $\{b\le X<b'\}$, $(b'-X)^+\le b'-b$, which contributes at most $b'-b$ to the ratio. On $\{X<b\}$, $e^{\theta X}\le e^{\theta b}$ for $\theta>0$,
while $\mathbb Ee^{\theta X}\ge e^{\theta b'}P(X>b')$ with $P(X>b')>0$; that part of the ratio is at most $e^{-\theta(b'-b)}\mathbb E|b'-X|/P(X>b')\to0$. So
$\liminf_{\theta\to\infty}\Lambda'(\theta)\ge b$, for every $b<M$.
\end{proof}

The \emph{rate function} is the Legendre transform
\[
\Lambda^*(x)=\sup_{\theta\in\R}\bigl(\theta x-\Lambda(\theta)\bigr)\in[0,\infty].
\]

\begin{lemma}[the rate function]\label{lem:ratefn}
\begin{enumerate}[label=(\roman*)]
\item $\Lambda^*$ is convex and lower semicontinuous, $\Lambda^*\ge0$ and $\Lambda^*(m)=0$.
\item If $x\ge m$, $\Lambda^*(x)=\sup_{\theta\ge0}(\theta x-\Lambda(\theta))$; if $x\le m$, the supremum may be taken over $\theta\le0$.
\item For $x\in(L,M)$, let $\theta_x$ be the solution of $\Lambda'(\theta_x)=x$. Then $\Lambda^*(x)=\theta_xx-\Lambda(\theta_x)$, $\Lambda^*$ is $C^1$ on $(L,M)$ with
$(\Lambda^*)'(x)=\theta_x$; it decreases on $(L,m]$ and increases on $[m,M)$.
\item $\Lambda^*(x)=\infty$ for $x\notin[L,M]$.
\end{enumerate}
\end{lemma}
\begin{proof}
(i) A supremum of affine functions of $x$ is convex and lower semicontinuous; $\theta=0$ gives $\Lambda^*\ge0$. By Jensen, $\Lambda(\theta)\ge\theta m$, so $\theta m-\Lambda(\theta)\le0$ and $\Lambda^*(m)=0$.
(ii) If $x\ge m$ and $\theta<0$, then $\theta x-\Lambda(\theta)\le\theta x-\theta m\le0$, which does not beat $\theta=0$. (iii) The function $\theta\mapsto\theta x-\Lambda(\theta)$ is concave with
derivative $x-\Lambda'(\theta)$, which vanishes at $\theta_x$, so the maximum is there. By the inverse function theorem $x\mapsto\theta_x$ is $C^1$, and differentiating
$\theta_xx-\Lambda(\theta_x)$ gives $\theta_x+(x-\Lambda'(\theta_x))\theta_x'=\theta_x$, negative for $x<m=\Lambda'(0)$ and positive for $x>m$. (iv) If $x>M$, then $\Lambda(\theta)\le\theta M$ for $\theta>0$
and $\theta x-\Lambda(\theta)\ge\theta(x-M)\to\infty$; similarly below $L$.
\end{proof}

Four rate functions to keep in mind (Exercise~\ref{ex:ratefns}): $\frac{x^2}2$ for $N(0,1)$; $x\log\frac x\lambda-x+\lambda$ for Poisson$(\lambda)$; $x-1-\log x$ for the exponential law
of mean $1$; and $I_p(x)$ for Bernoulli$(p)$. Figure~\ref{fig:ratefns} compares three of them with the same mean and variance: they agree near the mean,
where the central limit theorem rules, and nowhere else.

\begin{figure}[ht]
\centering\fig{ch07-ratefns}
\caption{Rate functions of three laws with mean $1$ and variance $1$: Gaussian $\frac{(x-1)^2}2$, Poisson$(1)$ $x\log x-x+1$, exponential $x-1-\log x$. All have curvature $1$
at the mean, as the central limit theorem demands, but their tails differ: the exponential law makes large values cheap, the Poisson law makes them expensive.}
\label{fig:ratefns}
\end{figure}

\begin{theorem}[Cram\'er]\label{thm:cramerR}
Let $X_1,X_2,\dots$ be i.i.d.\ copies of $X$, satisfying~\eqref{eq:cramercond}, and $S_n=X_1+\dots+X_n$.
\begin{enumerate}[label=(\alph*)]
\item For every $n\ge1$: $P(S_n\ge nx)\le e^{-n\Lambda^*(x)}$ if $x\ge m$, and $P(S_n\le nx)\le e^{-n\Lambda^*(x)}$ if $x\le m$.
\item For every closed $F\subseteq\R$, $\displaystyle\limsup_{n\to\infty}\frac1n\log P\Bigl(\frac{S_n}n\in F\Bigr)\le-\inf_F\Lambda^*$.
\item For every open $G\subseteq\R$, $\displaystyle\liminf_{n\to\infty}\frac1n\log P\Bigl(\frac{S_n}n\in G\Bigr)\ge-\inf_G\Lambda^*$.
\end{enumerate}
In particular $\frac1n\log P(S_n\ge na)\to-\Lambda^*(a)$ for $m<a<M$.
\end{theorem}
\begin{proof}
(a) For $\theta\ge0$, Markov's inequality gives $P(S_n\ge nx)\le e^{-n\theta x}\mathbb Ee^{\theta S_n}=e^{-n(\theta x-\Lambda(\theta))}$. Take the infimum over $\theta\ge0$ and use
Lemma~\ref{lem:ratefn}(ii). The lower tail is the same with $\theta\le0$.

(b) If $m\in F$ there is nothing to prove. Otherwise let $x_+$ be the smallest point of $F$ above $m$ and $x_-$ the largest below (if they exist; $F$ is closed). Then
$P(S_n/n\in F)\le P(S_n\ge nx_+)+P(S_n\le nx_-)\le2e^{-n\min(\Lambda^*(x_+),\Lambda^*(x_-))}$ by (a), and this minimum is $\inf_F\Lambda^*$, because $\Lambda^*$, convex with minimum
$0$ at $m$, is non-decreasing on $[m,\infty)$ and non-increasing on $(-\infty,m]$.

(c) It suffices to show that for every $x\in G$ with $\Lambda^*(x)<\infty$ and every small $\delta>0$,
\begin{equation}\label{eq:locallower}
\liminf_{n\to\infty}\frac1n\log P\Bigl(\Bigl|\frac{S_n}n-x\Bigr|<\delta\Bigr)\ge-\Lambda^*(x).
\end{equation}
First let $x\in(L,M)$ and $\theta=\theta_x$. Under the product of tilted laws $\mu_\theta^{\otimes n}$, the density of $\mu^{\otimes n}$ is $e^{-\theta(x_1+\dots+x_n)+n\Lambda(\theta)}$. On the event
$|S_n/n-x|<\delta$ this density is at least $e^{-n(\theta x-\Lambda(\theta))-n|\theta|\delta}$, so
\[
P\Bigl(\Bigl|\frac{S_n}n-x\Bigr|<\delta\Bigr)\ge e^{-n\Lambda^*(x)-n|\theta|\delta}\,P_\theta\Bigl(\Bigl|\frac{S_n}n-x\Bigr|<\delta\Bigr),
\]
where under $P_\theta$ the $X_i$ are i.i.d.\ with law $\mu_\theta$, of mean $x$ and finite variance. By the weak law the last probability tends to $1$, so the $\liminf$ in~\eqref{eq:locallower}
is at least $-\Lambda^*(x)-|\theta|\delta$. The probability decreases when $\delta$ decreases, so the same bound holds with any $\delta'<\delta$ in place of $\delta$ on the right; let
$\delta'\to0$.

If $x\notin(L,M)$ but $\Lambda^*(x)<\infty$, then $x$ is a finite endpoint, say $x=M$. For $0\le t<1$ the point $x_t=tM+(1-t)m$ lies in $(m,M)$, and by convexity
$\Lambda^*(x_t)\le t\Lambda^*(M)\le\Lambda^*(M)$. For $t$ close to $1$, $(x_t-\delta',x_t+\delta')\subseteq(x-\delta,x+\delta)$ for a small $\delta'$, and the first case applies at $x_t$.

Finally, for $m<a<M$, apply (b) to $[a,\infty)$ and (c) to $(a,\infty)$; both infima equal $\Lambda^*(a)$, by monotonicity and continuity of $\Lambda^*$ on $(L,M)$.
\end{proof}

The proof of (c) is the heart of the subject. To see a rare event, \emph{change the law} so that it becomes typical, here by tilting until the mean is $x$; the cost of the
change is the exponential of the relative entropy between the two laws, which is exactly $n\Lambda^*(x)$ (Exercise~\ref{ex:tiltentropy}). This gives an
algorithm, \emph{importance sampling}, which estimates probabilities of order $10^{-15}$ with a few thousand samples.

How good is the exponent alone? Cram\'er's 1938 paper already contained a sharper expansion; in the following form the result is due to Bahadur and Ranga Rao
(1960): if $X$ is not lattice,
\[
P(S_n\ge na)\sim\frac{e^{-n\Lambda^*(a)}}{\theta_a\sigma_a\sqrt{2\pi n}},\qquad\sigma_a^2=\Lambda''(\theta_a).
\]
For the Gaussian this is the tail estimate~\eqref{eq:gausstail} (Exercise~\ref{ex:gaussprefactor}). For an insurer facing $100$ independent claims with exponential law of mean $1$,
the probability that the total exceeds $200$ is $1.84\times10^{-15}$ (it is a gamma tail, computable exactly). The bound $e^{-n\Lambda^*(2)}$ gives $4.7\times10^{-14}$, the
Bahadur--Rao formula $1.88\times10^{-15}$, and the central limit theorem, $1-\Phi(10)$, gives $7.6\times10^{-24}$: wrong by a factor of $10^8$, and on the dangerous side.
That is the error Cram\'er warned the actuaries about.

\begin{think}
The Gaussian approximation is excellent at $a=m+\frac c{\sqrt n}$ and terrible at $a=m+c$. Where does it stop being good? Expand $\Lambda^*$ around $m$:
$\Lambda^*(m+h)=\frac{h^2}{2\sigma^2}-\frac{\kappa_3h^3}{6\sigma^6}+\dots$, with $\kappa_3$ the third cumulant. For which $h$, as a function of $n$, does the cubic term change $n\Lambda^*$ by
a bounded amount?
\end{think}

\section{Types and Sanov's theorem}

Now let $X$ take values in a finite alphabet $A$, with law $\mu$ such that $\mu(a)>0$ for every $a$. Write $\mathcal P(A)$ for the simplex of probability vectors on $A$, a
compact subset of $\R^A$, and $D(\nu\,\|\,\mu)=\sum_a\nu(a)\log\frac{\nu(a)}{\mu(a)}$, a continuous function of $\nu\in\mathcal P(A)$ since $\mu>0$. The \emph{type} of a word
$x=(x_1,\dots,x_n)\in A^n$ is its histogram $L_x(a)=\frac1n\#\{i:x_i=a\}$, and the random type of $X_1,\dots,X_n$ is the \emph{empirical measure} $L_n$. Let $\mathcal P_n$ be the set of
possible types of words of length $n$, and $T(\nu)=\{x\in A^n:L_x=\nu\}$ the \emph{type class} of $\nu\in\mathcal P_n$. Finally $H(\nu)=-\sum_a\nu(a)\log\nu(a)$.

\begin{lemma}[the method of types]\label{lem:types}
For $\nu\in\mathcal P_n$:
\begin{enumerate}[label=(\roman*)]
\item $|\mathcal P_n|\le(n+1)^{|A|}$;
\item every $x\in T(\nu)$ has $\mu^{\otimes n}(x)=e^{-n(H(\nu)+D(\nu\|\mu))}$;
\item $(n+1)^{-|A|}e^{nH(\nu)}\le|T(\nu)|\le e^{nH(\nu)}$;
\item $(n+1)^{-|A|}e^{-nD(\nu\|\mu)}\le P(L_n=\nu)\le e^{-nD(\nu\|\mu)}$.
\end{enumerate}
\end{lemma}
\begin{proof}
(i) Each $n\nu(a)$ is an integer in $\{0,\dots,n\}$. (ii)
\[
\mu^{\otimes n}(x)=\prod_a\mu(a)^{n\nu(a)}=\exp\Bigl(n\sum_a\nu(a)\log\mu(a)\Bigr),
\]
and $\sum_a\nu\log\mu=-H(\nu)-D(\nu\|\mu)$.
(iii) By (ii) with $\mu=\nu$ (allowing zeros, with $0^0=1$), every word of $T(\nu)$ has probability $e^{-nH(\nu)}$ under $\nu^{\otimes n}$, so $|T(\nu)|e^{-nH(\nu)}=\nu^{\otimes n}(T(\nu))\le1$.
For the lower bound we show that $T(\nu)$ is the most likely type class under $\nu^{\otimes n}$. For $\nu'\in\mathcal P_n$, counting words,
\[
\frac{\nu^{\otimes n}(T(\nu'))}{\nu^{\otimes n}(T(\nu))}=\prod_a\frac{(n\nu(a))!}{(n\nu'(a))!}\,\nu(a)^{n\nu'(a)-n\nu(a)}.
\]
If $\nu(a)=0<\nu'(a)$ the product is $0$. Otherwise use $\frac{k!}{l!}\le k^{k-l}$ for integers $k\ge1$, $l\ge0$ (for $l\le k$ it is a product of $k-l$ factors at most $k$; for
$l>k$ it is $1$ over a product of $l-k$ factors at least $k+1$); the factors with $\nu(a)=\nu'(a)=0$ equal $1$. Then the product is at most
$\prod_a(n\nu(a))^{n\nu(a)-n\nu'(a)}\nu(a)^{n\nu'(a)-n\nu(a)}=n^{\,n\sum_a(\nu(a)-\nu'(a))}=1$. Summing over the at most $(n+1)^{|A|}$ type classes,
$1\le(n+1)^{|A|}\nu^{\otimes n}(T(\nu))=(n+1)^{|A|}|T(\nu)|e^{-nH(\nu)}$. (iv) Multiply (ii) and (iii).
\end{proof}

So the probability of seeing the histogram $\nu$ is $e^{-nD(\nu\|\mu)}$, up to a polynomial factor that is invisible on the exponential scale. Boltzmann had made the same count in
1877 for the molecules of a gas: the number of microscopic states compatible with a histogram is $e^{nH}$.

\begin{theorem}[Sanov, finite alphabet]\label{thm:sanov}
For every $\Gamma\subseteq\mathcal P(A)$, with interior $\Gamma^\circ$ and closure $\overline\Gamma$ taken in $\mathcal P(A)$,
\[
-\inf_{\Gamma^\circ}D(\cdot\,\|\,\mu)\le\liminf_{n\to\infty}\frac1n\log P(L_n\in\Gamma)\le\limsup_{n\to\infty}\frac1n\log P(L_n\in\Gamma)\le-\inf_{\overline\Gamma}D(\cdot\,\|\,\mu).
\]
\end{theorem}
\begin{proof}
Upper bound: by Lemma~\ref{lem:types}, $P(L_n\in\Gamma)=\sum_{\nu\in\Gamma\cap\mathcal P_n}P(L_n=\nu)\le(n+1)^{|A|}e^{-n\inf_\Gamma D}$, and $\inf_\Gamma D\ge\inf_{\overline\Gamma}D$.
Lower bound: let $\nu\in\Gamma^\circ$. Fix $a_0\in A$, and let $\nu_n(a)=\lfloor n\nu(a)\rfloor/n$ for $a\ne a_0$ and $\nu_n(a_0)=1-\sum_{a\ne a_0}\nu_n(a)$. Then $\nu_n\in\mathcal P_n$ and
$|\nu_n(a)-\nu(a)|\le\frac{|A|}n$ for all $a$, so $\nu_n\in\Gamma$ for large $n$, and $P(L_n\in\Gamma)\ge P(L_n=\nu_n)\ge(n+1)^{-|A|}e^{-nD(\nu_n\|\mu)}$. By continuity of $D$,
$\liminf\frac1n\log P(L_n\in\Gamma)\ge-D(\nu\,\|\,\mu)$; take the supremum over $\nu\in\Gamma^\circ$.
\end{proof}

Sanov's theorem contains Cram\'er's for functions of $X$: if $f:A\to\R$, then $\frac1n\sum f(X_i)=\sum_af(a)L_n(a)$ is a linear function of $L_n$, and the exponent is a
minimum of relative entropy under a linear constraint. Computing it is Gibbs' variational principle.

\begin{proposition}[Gibbs' variational principle]\label{prop:gibbsvar}
Let $f:A\to\R$, $\langle f,\nu\rangle=\sum_af(a)\nu(a)$, $\Lambda(\theta)=\log\langle e^{\theta f},\mu\rangle$, and $\mu_\theta(a)=\mu(a)e^{\theta f(a)-\Lambda(\theta)}$. For every $\nu\in\mathcal P(A)$ and $\theta\in\R$,
\[
D(\nu\,\|\,\mu)=\theta\langle f,\nu\rangle-\Lambda(\theta)+D(\nu\,\|\,\mu_\theta).
\]
Consequently, if $f$ is not constant and $\langle f,\mu\rangle\le x<\max f$, then $\min\{D(\nu\,\|\,\mu):\langle f,\nu\rangle\ge x\}=\Lambda^*(x)$, attained only at $\nu=\mu_{\theta_x}$.
\end{proposition}
\begin{proof}
$\log\frac{\nu(a)}{\mu(a)}=\log\frac{\nu(a)}{\mu_\theta(a)}+\theta f(a)-\Lambda(\theta)$; average over $\nu$. Here $\Lambda$ is the cumulant generating function of $f(X)$, whose essential
supremum is $\max f$. If $\langle f,\nu\rangle=y\ge x$, take $\theta=\theta_x\ge0$: $D(\nu\|\mu)=\theta_xy-\Lambda(\theta_x)+D(\nu\|\mu_{\theta_x})\ge\theta_xx-\Lambda(\theta_x)=\Lambda^*(x)$, with
equality only if $\theta_x(y-x)=0$ and $\nu=\mu_{\theta_x}$; and $\nu=\mu_{\theta_x}$ does satisfy $\langle f,\nu\rangle=\Lambda'(\theta_x)=x$.
\end{proof}

The minimizer is a \emph{Gibbs measure}: among all histograms with the required average, the cheapest one tilts $\mu$ exponentially. That is why Boltzmann's distribution
$e^{-E/kT}$ appears in physics, and the next theorem makes the statement precise.

\section{The shape of a rare event}

\begin{theorem}[Gibbs conditioning principle]\label{thm:gibbscond}
Let $A$, $\mu$, $f$ be as above, $X_1,X_2,\dots$ i.i.d.\ with law $\mu$, $\langle f,\mu\rangle<x<\max f$, $E_n=\{f(X_1)+\dots+f(X_n)\ge nx\}$, and $\nu^*=\mu_{\theta_x}$. For each fixed $k$
and $a_1,\dots,a_k\in A$,
\[
P(X_1=a_1,\dots,X_k=a_k\mid E_n)\longrightarrow\nu^*(a_1)\cdots\nu^*(a_k).
\]
\end{theorem}
\begin{proof}
\emph{The empirical measure concentrates.} Let $\Gamma=\{\nu:\langle f,\nu\rangle\ge x\}$, closed, so $E_n=\{L_n\in\Gamma\}$. Fix $\delta>0$ and let $K=\{\nu\in\Gamma:\max_a|\nu(a)-\nu^*(a)|\ge\delta\}$, a
compact set on which the continuous function $D(\cdot\|\mu)$ has a minimum, larger than $\Lambda^*(x)$ by Proposition~\ref{prop:gibbsvar}: say $\Lambda^*(x)+2\eta$. By Theorem~\ref{thm:sanov},
$P(L_n\in K)\le e^{-n(\Lambda^*(x)+\eta)}$ for large $n$. On the other side, $\Gamma^\circ$ contains $\{\nu:\langle f,\nu\rangle>x\}$, which contains $\mu_{\theta_y}$ for $x<y<\max f$, with
$D(\mu_{\theta_y}\|\mu)=\Lambda^*(y)\to\Lambda^*(x)$ as $y\downarrow x$. So $\inf_{\Gamma^\circ}D=\Lambda^*(x)$, and $P(E_n)\ge e^{-n(\Lambda^*(x)+\eta/2)}$ for large $n$. Therefore
$P(L_n\in K\mid E_n)\to0$.

\emph{Given its type, a word is uniform.} By Lemma~\ref{lem:types}(ii), all words with type $\nu$ are equally likely, so given $L_n=\nu$ the first $k$ letters are $k$ draws without
replacement from an urn with $n\nu(a)$ balls of colour $a$:
\[
P(X_1=a_1,\dots,X_k=a_k\mid L_n=\nu)=\prod_{i=1}^k\frac{(n\nu(a_i)-c_i)^+}{n-i+1},
\]
where $c_i\le i-1$ is the number of $j<i$ with $a_j=a_i$ (the positive part makes the formula right when the urn runs out of colour $a_i$). Each factor is in $[0,1]$
and differs from $\nu(a_i)$ by at most $\frac{|c_i-(i-1)\nu(a_i)|}{n-i+1}\le\frac{i-1}{n-k}$, and a product of numbers in $[0,1]$ moves by at most the sum of the changes, so
$\bigl|P(X_1=a_1,\dots,X_k=a_k\mid L_n=\nu)-\prod_i\nu(a_i)\bigr|\le\frac{k^2}{n-k}$.

\emph{Conclusion.} $E_n$ depends only on $L_n$, so averaging over $L_n$ given $E_n$: the conditional probability is within $\frac{k^2}{n-k}+P(L_n\in K\mid E_n)$ of a product
$\prod_i\nu(a_i)$ with $\max_a|\nu(a)-\nu^*(a)|<\delta$, hence within $\frac{k^2}{n-k}+P(L_n\in K\mid E_n)+k\delta$ of $\prod_i\nu^*(a_i)$. Let $n\to\infty$, then $\delta\to0$.
\end{proof}

\begin{example}[the Brandeis dice]\label{ex:brandeis}
At a summer school at Brandeis University in 1962, Edwin Jaynes posed a problem: a die has been tossed a very large number of times, and the average number of
spots was $4.5$ rather than $3.5$. What should we predict for the next toss? Jaynes answered with the maximum entropy principle; Theorem~\ref{thm:gibbscond} gives the
same answer for a fair die conditioned on a rare average. With $f(a)=a$ on $\{1,\dots,6\}$ and $\mu$ uniform, solve $\Lambda'(\theta)=4.5$: $\theta=0.371$ and
\[
\nu^*=(0.054,\ 0.079,\ 0.114,\ 0.165,\ 0.240,\ 0.348)\quad\text{on }1,\dots,6.
\]
Each die in the lucky sequence behaves as if its probabilities were tilted geometrically: not ``a few sixes and the rest normal'', but a little more of everything above
$3.5$. The rate is $\Lambda^*(4.5)=0.178$; for $100$ dice, $P(\text{sum}\ge450)=1.5\times10^{-9}$ exactly, against the upper bound $e^{-17.8}=1.8\times10^{-8}$.
\end{example}

The conditioning principle explains why the heavy-tailed sums of Exercise~\ref{ex:heavyld} behave so differently: there the cheapest way to a large sum is one huge
term, not a uniform tilt.

\section{The best test}

A coin is either fair ($\mu_0$) or biased with $P(\text{heads})=0.6$ ($\mu_1$). We see $n$ tosses and must decide. A test is a set $B\subseteq A^n$ of outcomes on which we
declare ``fair''. It errs in two ways: rejecting a fair coin, with probability $\mu_0^{\otimes n}(B^c)$, and accepting a biased one, with probability $\mu_1^{\otimes n}(B)$. A spam filter,
a medical screening and an A/B test all face the same trade-off. Let
\[
\beta_n(\varepsilon)=\min\bigl\{\mu_1^{\otimes n}(B):\ B\subseteq A^n,\ \mu_0^{\otimes n}(B^c)\le\varepsilon\bigr\}
\]
be the smallest second error compatible with a first error at most $\varepsilon$.

\begin{theorem}[Chernoff--Stein lemma]\label{thm:chernoffstein}
Let $\mu_0,\mu_1$ be laws on a finite set $A$ with full support. For every $0<\varepsilon<1$, $\displaystyle\frac1n\log\beta_n(\varepsilon)\to-D(\mu_0\,\|\,\mu_1)$.
\end{theorem}
\begin{proof}
Write $D=D(\mu_0\|\mu_1)$ and $P_j=\mu_j^{\otimes n}$.

\emph{A good test.} Fix $\delta>0$ and let $B=\{x:\max_a|L_x(a)-\mu_0(a)|<\delta\}$. By the weak law applied to each indicator $\mathbf 1\{X_i=a\}$, $P_0(B^c)\to0$, so $P_0(B^c)\le\varepsilon$ for
large $n$. And by Lemma~\ref{lem:types}, $P_1(B)\le(n+1)^{|A|}e^{-nd_\delta}$, where $d_\delta=\min\{D(\nu\|\mu_1):\max_a|\nu(a)-\mu_0(a)|\le\delta\}\to D$ as $\delta\to0$, by continuity. So
$\limsup\frac1n\log\beta_n(\varepsilon)\le-d_\delta$ for every $\delta$.

\emph{No test is better.} The log-likelihood ratio $\frac1n\log\frac{P_0(X)}{P_1(X)}=\frac1n\sum_i\log\frac{\mu_0(X_i)}{\mu_1(X_i)}$ is an average of bounded i.i.d.\ variables of mean $D$ under $P_0$.
By the weak law, $T=\{x:P_1(x)\ge P_0(x)e^{-n(D+\delta)}\}$ has $P_0(T)\to1$. If $P_0(B^c)\le\varepsilon$, then $P_0(B\cap T)\ge1-\varepsilon-P_0(T^c)\ge\frac{1-\varepsilon}2$ for large $n$ (uniformly in $B$), and
\[
P_1(B)\ge\sum_{x\in B\cap T}P_1(x)\ge e^{-n(D+\delta)}P_0(B\cap T)\ge\frac{1-\varepsilon}2e^{-n(D+\delta)}.
\]
So $\liminf\frac1n\log\beta_n(\varepsilon)\ge-D-\delta$.
\end{proof}

The lemma carries the names of Herman Chernoff and of the Charles Stein of the previous chapter. It says that relative entropy is not only a measure of surprise but the
exact price of telling two hypotheses apart, and that the first error $\varepsilon$ does not affect the exponent. For the coin, $D(\mu_0\|\mu_1)=0.0204$, so pushing the second error
down to $10^{-6}$ costs about $\log(10^6)/0.0204\approx680$ tosses, to leading order; the exact optimum for $\varepsilon=0.05$, computed numerically, is about $1000$, because
the correction to the exponent is of order $n^{-1/2}$. Note the asymmetry: $D(\mu_1\|\mu_0)=0.0201$ is the exponent if the roles are exchanged.

\section{Varadhan's lemma and a toy magnet}

Laplace's method says that $\int e^{nF(x)}g(x)\dd x$ grows like $e^{n\max F}$. Varadhan's lemma is its probabilistic form: an exponential integral against a law satisfying
large deviation bounds is governed by a balance between the integrand and the rate.

\begin{theorem}[Varadhan's lemma, for sample means]\label{thm:varadhan}
Under the assumptions of Theorem~\ref{thm:cramerR}, let $F:\R\to\R$ be continuous and bounded. Then
\[
\lim_{n\to\infty}\frac1n\log\mathbb E\,e^{nF(S_n/n)}=\sup_{x\in\R}\bigl(F(x)-\Lambda^*(x)\bigr),
\]
and for every closed $C\subseteq\R$,
\[
\limsup_{n\to\infty}\frac1n\log\mathbb E\Bigl[e^{nF(S_n/n)};\ \frac{S_n}n\in C\Bigr]\le\sup_C(F-\Lambda^*).
\]
\end{theorem}
\begin{proof}
Let $s=\sup(F-\Lambda^*)$; it is finite, between $F(m)$ and $\sup F$.

\emph{Lower bound.} Let $x$ with $\Lambda^*(x)<\infty$ and $\varepsilon>0$; choose $\delta$ with $F\ge F(x)-\varepsilon$ on $(x-\delta,x+\delta)$. Then
$\mathbb Ee^{nF(S_n/n)}\ge e^{n(F(x)-\varepsilon)}P(|S_n/n-x|<\delta)$, and Theorem~\ref{thm:cramerR}(c) gives $\liminf\frac1n\log\mathbb Ee^{nF(S_n/n)}\ge F(x)-\varepsilon-\Lambda^*(x)$.

\emph{Upper bound on a closed set $C$.} Taking $\theta=\pm1$ in the definition, $\Lambda^*(y)\ge|y|-c$ with $c=\max(\Lambda(1),\Lambda(-1))$. Fix $K>0$ and choose $R$ with $R-c>\sup|F|+K+1$.
By Theorem~\ref{thm:cramerR}(b) applied to $\{|y|\ge R\}$, $\mathbb E[e^{nF(S_n/n)};|S_n/n|\ge R]\le e^{n\sup|F|}e^{-n(R-c-1)}\le e^{-nK}$ for large $n$. Next, $F$ is uniformly
continuous on $[-R,R]$: cover $C\cap[-R,R]$ by finitely many closed intervals $J_1,\dots,J_k$, of length so small that $F$ varies by at most $\varepsilon$ on each. By Theorem~\ref{thm:cramerR}(b) for the
closed set $C\cap J_i$,
\[
\limsup_{n\to\infty}\frac1n\log\mathbb E\Bigl[e^{nF(S_n/n)};\frac{S_n}n\in C\cap J_i\Bigr]\le\sup_{C\cap J_i}F-\inf_{C\cap J_i}\Lambda^*\le\sup_{C\cap J_i}(F-\Lambda^*)+\varepsilon.
\]
The $\limsup$ of $\frac1n\log$ of a sum of $k+1$ terms is the largest of the $\limsup$s, so the left side of the claim is at most $\max\bigl(\sup_C(F-\Lambda^*)+\varepsilon,-K\bigr)$. Let
$\varepsilon\to0$ and $K\to\infty$. With $C=\R$ this is the upper bound for the first statement.
\end{proof}

\emph{The Curie--Weiss magnet.} In his thesis of 1895 Pierre Curie found that iron loses its magnetism above a critical temperature, now the Curie point. In 1907 Pierre
Weiss explained it with a ``molecular field'': each elementary magnet feels the average of all the others. The simplest model has $n$ spins $\sigma_i\in\{-1,+1\}$, an energy
$-\frac1{2n}(\sigma_1+\dots+\sigma_n)^2$ that rewards alignment, and, at inverse temperature $\beta>0$, the Gibbs probability
\[
P_{n,\beta}(\sigma)=\frac1{Z_n}\exp\Bigl(\frac\beta{2n}\Bigl(\sum_i\sigma_i\Bigr)^2\Bigr),\qquad Z_n=\sum_{\sigma\in\{\pm1\}^n}\exp\Bigl(\frac\beta{2n}\Bigl(\sum_i\sigma_i\Bigr)^2\Bigr).
\]
The \emph{magnetization} is $m_n=\frac1n\sum_i\sigma_i$. If the $\sigma_i$ are fair independent signs, $Z_n=2^n\mathbb E\,e^{nF(m_n)}$ with $F(y)=\frac\beta2\min(y^2,1)$, which is bounded
and continuous and equals $\frac\beta2y^2$ on $[-1,1]$, where $m_n$ lives. For a fair sign, $\Lambda(\theta)=\log\cosh\theta$ and
\[
\Lambda^*(y)=J(y)=\frac{1+y}2\log(1+y)+\frac{1-y}2\log(1-y)\quad\text{on }[-1,1]
\]
(Exercise~\ref{ex:ratefns}).

\begin{theorem}[phase transition]\label{thm:curieweiss}
$\frac1n\log Z_n\to\log2+\max_{|y|\le1}\bigl(\frac\beta2y^2-J(y)\bigr)$. If $\beta\le1$ the maximum is attained only at $y=0$; if $\beta>1$ it is attained exactly at $\pm m^*(\beta)$, where
$m^*(\beta)$ is the unique positive solution of $m=\tanh(\beta m)$. In both cases, for every $\delta>0$, $P_{n,\beta}\bigl(\operatorname{dist}(m_n,\mathcal M)\ge\delta\bigr)\to0$ exponentially fast,
where $\mathcal M$ is the set of maximizers.
\end{theorem}
\begin{proof}
The limit is Theorem~\ref{thm:varadhan}, since $\Lambda^*=\infty$ outside $[-1,1]$. Let $g(y)=\frac\beta2y^2-J(y)$, even, continuous on $[-1,1]$, with $g'(y)=\beta y-\operatorname{artanh}y$ on
$(-1,1)$ and $g'(y)\to-\infty$ as $y\to1$; so maximizers are interior critical points, solutions of $y=\tanh(\beta y)$. If $\beta\le1$, $\tanh(\beta y)<\beta y\le y$ for $y>0$, and $0$ is the only
critical point. If $\beta>1$, $h(y)=\tanh(\beta y)-y$ is strictly concave on $[0,\infty)$, with $h(0)=0$, $h'(0)=\beta-1>0$ and $h(1)<0$: it has exactly one zero $m^*$ in $(0,1)$. Then
$g'>0$ on $(0,m^*)$ and $g'<0$ on $(m^*,1)$, so $m^*$ is the maximum on $[0,1]$, and by evenness $\pm m^*$ are the maximizers, while $0$ is not one.

For the concentration, let $C=\{y\in[-1,1]:\operatorname{dist}(y,\mathcal M)\ge\delta\}$, compact. On $C$ the continuous $g$ has maximum $<\max g$, say $\max g-2\eta$. Then
\[
P_{n,\beta}\bigl(m_n\in C\bigr)=\frac{\mathbb E[e^{nF(m_n)};m_n\in C]}{\mathbb Ee^{nF(m_n)}},
\]
whose numerator is at most $e^{n(\max g-\eta)}$ for large $n$ by the second part of Theorem~\ref{thm:varadhan}, and whose denominator is at least $e^{n(\max g-\eta/2)}$.
\end{proof}

At high temperature ($\beta<1$) the spins disagree and the magnet is not magnetized. Below the Curie point ($\beta>1$) the magnetization is close to $+m^*$ or $-m^*$, each with
probability $\frac12$ by symmetry: $m^*(1.5)=0.859$, $m^*(2)=0.958$. The system must choose a direction, and the symmetric average hides it. Figure~\ref{fig:curieweiss} shows the
function $g$ changing shape at $\beta=1$. The square-root law $m^*(\beta)\sim\sqrt{3(\beta-1)}$ near the transition is Exercise~\ref{ex:curieweiss}. At $\beta=1$ exactly, $g(y)=-\frac{y^4}{12}+O(y^6)$ is so flat that
$m_n$ lives at the scale $n^{-1/4}$ instead of $n^{-1/2}$: Ellis and Newman (1978) proved that $n^{1/4}m_n$ converges in law to the density proportional to $e^{-u^4/12}$, a
non-Gaussian limit for a sum of $\pm1$ spins.

\begin{figure}[ht]
\centering\fig{ch07-curieweiss}
\caption{The function $g(y)=\frac\beta2y^2-J(y)$ of the Curie--Weiss model, whose maximum gives the free energy, for $\beta=0.8$, $1$ and $1.5$. At $\beta=1$ the maximum at $0$
becomes flat; beyond it splits into two, at $\pm m^*(\beta)$ (dots), and the magnet magnetizes.}
\label{fig:curieweiss}
\end{figure}

\section{Exercises}

\begin{exercise}[$\star$]\label{ex:ratefns}
Compute $\Lambda^*$ for $N(0,1)$, Poisson$(\lambda)$, the exponential law of mean $1$ (for which $\Lambda(\theta)=\infty$ when $\theta\ge1$; check that Theorem~\ref{thm:cramerR}(a) still
holds), Bernoulli$(p)$, and a uniform sign $\pm1$.
\end{exercise}

\begin{exercise}[$\star$]\label{ex:tiltentropy}
Show that $D(\mu_\theta\,\|\,\mu)=\theta\Lambda'(\theta)-\Lambda(\theta)=\Lambda^*(\Lambda'(\theta))$ (for a law on a finite set, or using densities in general). So the cost of tilting the mean to $x$
is exactly the rate $\Lambda^*(x)$.
\end{exercise}

\begin{exercise}[$\star$]\label{ex:gaussprefactor}
For $N(0,1)$ variables, check the Bahadur--Rao formula: $P(S_n\ge na)\sim\frac{e^{-na^2/2}}{a\sqrt{2\pi n}}$ for $a>0$.
\end{exercise}

\begin{exercise}[$\star\star$]\label{ex:endpoint}
Suppose $M=\operatorname{ess\,sup}X<\infty$ and $P(X=M)=p>0$. Show that $\Lambda^*(M)=-\log p$ and $P(S_n\ge nM)=p^n$. What if $P(X=M)=0$?
\end{exercise}

\begin{exercise}[$\star\star$]\label{ex:heavyld}
Let $X\ge0$ with $P(X>t)=(1+t)^{-3}$. Show that $\Lambda(\theta)=\infty$ for $\theta>0$, that $\Lambda^*(x)=0$ for every $x\ge m$, and that for $a>m$,
$P(S_n\ge na)$ decays only polynomially: $P(S_n\ge na)\ge\frac12P\bigl(X_1\ge n(a-m)+m+n^{2/3}\bigr)$ for large $n$. How does the large sum happen?
\end{exercise}

\begin{exercise}[$\star\star$]\label{ex:testcoin}
In the setting of Theorem~\ref{thm:chernoffstein}, suppose instead that we want both errors small, with the same exponent. Show that the likelihood ratio test
$B=\{x:P_0(x)\ge P_1(x)\}$ has both errors at most $e^{-nC}$, where $C=-\min_{0\le s\le1}\log\sum_a\mu_0(a)^{1-s}\mu_1(a)^s$ (the \emph{Chernoff information}), and compute $C$ for the fair coin
against $P(\text{heads})=0.6$.
\end{exercise}

\begin{exercise}[$\star\star$]\label{ex:curieweiss}
In Theorem~\ref{thm:curieweiss}, show that $m^*(\beta)\sim\sqrt{3(\beta-1)}$ as $\beta\downarrow1$, and that for $\beta>1$ and small $\delta>0$, $P_{n,\beta}(|m_n-m^*|<\delta)\to\frac12$.
\end{exercise}

\section{Hints}

\begin{hint}{ex:ratefns}
Solve $\Lambda'(\theta)=x$. For the exponential law, $\Lambda(\theta)=-\log(1-\theta)$ for $\theta<1$; Markov's inequality only needs $\theta$ where $\Lambda$ is finite.
\end{hint}
\begin{hint}{ex:tiltentropy}
$\log\frac{\dd\mu_\theta}{\dd\mu}=\theta x-\Lambda(\theta)$; integrate against $\mu_\theta$.
\end{hint}
\begin{hint}{ex:gaussprefactor}
$S_n/\sqrt n$ is exactly $N(0,1)$; use~\eqref{eq:gausstail}.
\end{hint}
\begin{hint}{ex:endpoint}
$\Lambda(\theta)=\theta M+\log\mathbb Ee^{\theta(X-M)}$, and $\mathbb Ee^{\theta(X-M)}\downarrow P(X=M)$ as $\theta\to\infty$.
\end{hint}
\begin{hint}{ex:heavyld}
One large value and $n-1$ ordinary ones: by the weak law, $X_2+\dots+X_n\ge(n-1)m-n^{2/3}$ with probability tending to $1$.
\end{hint}
\begin{hint}{ex:testcoin}
For each $s\in[0,1]$, $\mathbf 1\{P_0(x)<P_1(x)\}\le P_0(x)^{-s}P_1(x)^s$. Multiply by $P_0(x)$ and sum.
\end{hint}
\begin{hint}{ex:curieweiss}
Expand $\tanh u=u-\frac{u^3}3+O(u^5)$, after checking that $m^*(\beta)\to0$. For the second part, use the symmetry $\sigma\mapsto-\sigma$.
\end{hint}

\section{Solutions}

\begin{solution}{ex:ratefns}
\idea Each rate function is $\theta_xx-\Lambda(\theta_x)$ at the tilt that moves the mean to $x$.
\writeup Gaussian: $\Lambda=\frac{\theta^2}2$, $\theta_x=x$, $\Lambda^*=\frac{x^2}2$. Poisson: $\Lambda=\lambda(e^\theta-1)$, $\theta_x=\log\frac x\lambda$ for $x>0$, $\Lambda^*=x\log\frac x\lambda-x+\lambda$; at
$x=0$, $\Lambda^*(0)=\sup_\theta(-\lambda(e^\theta-1))=\lambda$, which is the same formula with $0\log0=0$; $\Lambda^*=\infty$ for $x<0$. Exponential: $\Lambda=-\log(1-\theta)$ for $\theta<1$ and $\infty$
otherwise; $\theta_x=1-\frac1x$, $\Lambda^*=x-1-\log x$ for $x>0$. In the proof of Theorem~\ref{thm:cramerR}(a), Markov's inequality with $\theta\in[0,1)$ gives
$P(S_n\ge nx)\le e^{-n(\theta x-\Lambda(\theta))}$, and the supremum over these $\theta$ is still $\Lambda^*(x)$ for $x\ge1$. Bernoulli$(p)$: $\Lambda=\log(1-p+pe^\theta)$, and the
computation of Tome~3, Theorem~4.7 gives $I_p(x)$ on $[0,1]$. Uniform sign: $X=2B-1$ with $B$ Bernoulli$(\frac12)$, so $\Lambda(\theta)=\log\cosh\theta$ and
$\Lambda^*(y)=I_{1/2}(\frac{1+y}2)=\frac{1+y}2\log(1+y)+\frac{1-y}2\log(1-y)$, since $\frac{1+y}2\log\frac{(1+y)/2}{1/2}=\frac{1+y}2\log(1+y)$.
\end{solution}

\begin{solution}{ex:tiltentropy}
\idea Relative entropy is the mean log-density, and the log-density of a tilt is affine.
\writeup $D(\mu_\theta\|\mu)=\int\log\frac{\dd\mu_\theta}{\dd\mu}\dd\mu_\theta=\int(\theta x-\Lambda(\theta))\mu_\theta(\dd x)=\theta\Lambda'(\theta)-\Lambda(\theta)$, and by Lemma~\ref{lem:ratefn}(iii) with
$x=\Lambda'(\theta)$, $\theta=\theta_x$, this is $\Lambda^*(x)$. In the proof of Theorem~\ref{thm:cramerR}(c), the product of $n$ tilted laws costs $e^{-nD(\mu_\theta\|\mu)}$ on its typical
event: the price of making the rare event typical.
\end{solution}

\begin{solution}{ex:gaussprefactor}
\idea For Gaussians everything is explicit: $S_n\sim N(0,n)$.
\writeup $P(S_n\ge na)=P(Z\ge a\sqrt n)$. By~\eqref{eq:gausstail}, $\frac{x}{1+x^2}\varphi(x)\le P(Z\ge x)$, and the upper bound $P(Z\ge x)\le\frac{\varphi(x)}x$ follows from
$\int_x^\infty\varphi(t)\dd t\le\int_x^\infty\frac txe^{-t^2/2}\frac{\dd t}{\sqrt{2\pi}}$. With $x=a\sqrt n$, $P(S_n\ge na)\sim\frac{\varphi(a\sqrt n)}{a\sqrt n}=\frac{e^{-na^2/2}}{a\sqrt{2\pi n}}$. Here $\theta_a=a$,
$\sigma_a=1$, as in the general formula.
\end{solution}

\begin{solution}{ex:endpoint}
\idea At the top of the support there is nothing to tilt towards: the only way to average $M$ is to be $M$ every time.
\writeup $S_n\ge nM$ with $X_i\le M$ forces all $X_i=M$, so $P(S_n\ge nM)=p^n$. Next $\theta M-\Lambda(\theta)=-\log\mathbb Ee^{\theta(X-M)}$; as $\theta\to\infty$, $e^{\theta(X-M)}$ decreases to
$\mathbf 1\{X=M\}$, so by monotone convergence $-\log\mathbb Ee^{\theta(X-M)}$ increases to $-\log p$, and for $\theta\le0$ the expression is at most $0$ by Lemma~\ref{lem:ratefn}(ii). So
$\Lambda^*(M)=-\log p$, and the bound of Theorem~\ref{thm:cramerR}(a) is an equality for every $n$. If $P(X=M)=0$, then $P(S_n\ge nM)=0$ and $\Lambda^*(M)=\infty$, consistent with the convention $e^{-\infty}=0$.
\end{solution}

\begin{solution}{ex:heavyld}
\idea Without exponential moments, a large deviation is caused by one variable doing all the work, which costs only a power of $n$.
\writeup For $\theta>0$, $\mathbb Ee^{\theta X}\ge e^{\theta t}(1+t)^{-3}\to\infty$ as $t\to\infty$, so $\Lambda(\theta)=\infty$. For $\theta\le0$, $\theta x-\Lambda(\theta)\le\theta(x-m)\le0$ when $x\ge m$ (Jensen), so $\Lambda^*(x)=0$.
The variance is finite ($\mathbb EX^2=2\int_0^\infty t(1+t)^{-3}\dd t<\infty$), so by Chebyshev $P(X_2+\dots+X_n<(n-1)m-n^{2/3})\le\frac{(n-1)\operatorname{Var}X}{n^{4/3}}\to0$; let $n$ be so
large that this is at most $\frac12$. On the intersection of the independent events $\{X_1\ge n(a-m)+m+n^{2/3}\}$ and $\{X_2+\dots+X_n\ge(n-1)m-n^{2/3}\}$, the sum is
$S_n\ge n(a-m)+m+(n-1)m=na$. So $P(S_n\ge na)\ge\frac12P(X_1\ge n(a-m)+m+n^{2/3})=\frac12\bigl(1+n(a-m)+m+n^{2/3}\bigr)^{-3}$, which is of order $n^{-3}$. The sum is large because one claim is enormous, and the others are
typical: the opposite of the Gibbs picture.
\end{solution}

\begin{solution}{ex:testcoin}
\idea Bound an indicator by a power of the likelihood ratio; optimizing the power is a Chernoff bound, and the exponent is symmetric in the two errors.
\writeup On $B^c$, $P_0(x)<P_1(x)$, so for $s\in[0,1]$, $P_0(x)\le P_0(x)^{1-s}P_1(x)^s$, and $P_0(B^c)\le\sum_xP_0(x)^{1-s}P_1(x)^s=\bigl(\sum_a\mu_0(a)^{1-s}\mu_1(a)^s\bigr)^n$, since the sum
factorizes over the coordinates. On $B$, $P_1(x)\le P_1(x)^sP_0(x)^{1-s}$, the same bound. Minimize over $s$. For the coins, $\sum_a\mu_0(a)^{1-s}\mu_1(a)^s=0.5^{1-s}(0.6^s+0.4^s)$, minimal at $s=0.498$, which gives $C=0.00508$, about a quarter of $D(\mu_0\|\mu_1)=0.0204$.
Demanding that both errors be small costs four times as many tosses as the Chernoff--Stein regime, where one error is only required to stay below a fixed $\varepsilon$.
\end{solution}

\begin{solution}{ex:curieweiss}
\idea Near the transition, $\tanh$ is a line minus a cube, and the positive fixed point is where the two balance.
\writeup \emph{$m^*\to0$.} Let $h_\beta(y)=\tanh(\beta y)-y$; the proof of Theorem~\ref{thm:curieweiss} shows $h_\beta>0$ on $(0,m^*)$ and $h_\beta<0$ on $(m^*,1]$. Fix
$y_0\in(0,1)$. As $\beta\downarrow1$, $h_\beta(y_0)\to\tanh y_0-y_0<0$, so $h_\beta(y_0)<0$ and $m^*(\beta)<y_0$ for $\beta$ close to $1$.
\emph{The asymptotics.} Now divide $m^*=\tanh(\beta m^*)=\beta m^*-\frac{(\beta m^*)^3}3+O(m^{*5})$ by $m^*$: $1=\beta-\frac{\beta^3m^{*2}}3+O(m^{*4})$, so
$m^{*2}=\frac{3(\beta-1)}{\beta^3}(1+O(m^{*2}))$ and $m^*\sim\sqrt{3(\beta-1)}$.
\emph{The two states.} The Gibbs weight depends on $\sigma$ only through $(\sum\sigma_i)^2$, so $\sigma\mapsto-\sigma$ preserves $P_{n,\beta}$ and $m_n$ is symmetric. For $\delta<m^*$ the intervals
$(m^*-\delta,m^*+\delta)$ and its mirror image are disjoint, have equal probabilities, and by Theorem~\ref{thm:curieweiss} their union has probability tending to $1$. Each tends to $\frac12$.
\end{solution}

\section*{Sources and further reading}

\begin{itemize}
\item \emph{History.} F.~Esscher, ``On the probability function in the collective theory of risk'', \emph{Skand.\ Aktuarietidskr.}\ 15 (1932), 175--195. H.~Cram\'er, ``Sur un nouveau
th\'eor\`eme-limite de la th\'eorie des probabilit\'es'', in \emph{Colloque consacr\'e \`a la th\'eorie des probabilit\'es}, Actualit\'es scientifiques et industrielles 736, Hermann, 1938
(the Geneva conference of 1937); English translation by H.~Touchette, arXiv:1802.05988, whose introduction tells the story. I.~N.~Sanov, ``On the probability of large
deviations of random variables'' (in Russian), \emph{Mat.\ Sb.}\ 42(84) (1957), 11--44. S.~R.~S.~Varadhan, ``Asymptotic probabilities and differential equations'', \emph{Comm.\ Pure
Appl.\ Math.}\ 19 (1966), 261--286. S.~R.~S.~Varadhan, ``Large deviations'', \emph{Ann.\ Probab.}\ 36 (2008), 397--419 (his own overview, with history). Abel Prize committee, citation for
S.~R.~S.~Varadhan, 2007. R.~R.~Bahadur and R.~Ranga Rao, ``On deviations of the sample mean'', \emph{Ann.\ Math.\ Statist.}\ 31 (1960), 1015--1027.
\item E.~T.~Jaynes, ``Information theory and statistical mechanics'', in \emph{Statistical Physics}, Brandeis Summer Institute 1962, vol.~3, Benjamin, 1963, 181--218 (the dice).
P.~Curie, ``Propri\'et\'es magn\'etiques des corps \`a diverses temp\'eratures'', \emph{Ann.\ Chim.\ Phys.}, 1895 (his thesis). P.~Weiss, ``L'hypoth\`ese du champ mol\'eculaire et la propri\'et\'e
ferromagn\'etique'', \emph{J.\ Phys.\ Th\'eor.\ Appl.}\ 6 (1907), 661--690.
\item A.~Dembo and O.~Zeitouni, \emph{Large Deviations Techniques and Applications}, 2nd ed., Springer, 1998. The standard reference: Cram\'er without condition~\eqref{eq:cramercond},
Sanov on Polish spaces, the contraction principle, Varadhan's lemma in general, and the Gibbs conditioning principle.
\item T.~M.~Cover and J.~A.~Thomas, \emph{Elements of Information Theory}, 2nd ed., Wiley, 2006, Chapter~11. The method of types, Sanov, and the Chernoff--Stein lemma
from the information theorist's side.
\item R.~S.~Ellis, \emph{Entropy, Large Deviations, and Statistical Mechanics}, Springer, 1985. The Curie--Weiss model and its relatives, with full proofs. R.~S.~Ellis and
C.~M.~Newman, ``Limit theorems for sums of dependent random variables occurring in statistical mechanics'', \emph{Z.\ Wahrsch.\ Verw.\ Gebiete} 44 (1978), 117--139.
\item F.~den Hollander, \emph{Large Deviations}, Fields Institute Monographs 14, AMS, 2000. A short and lively introduction.
\end{itemize}
